\section{Unrestricted pro-\texorpdfstring{$3$}{3} Selmer cuspidality}
\subsection{Compact cancellation and the unique abelian cusp parameter}
\label{subsec:lm-unrestricted-reduction}

Let $n=3^r$ with $r\ge1$, put $J_n=\operatorname{Jac}(X_n)$, and write
$N_n=\pi_1(U_{n,\overline\Q})^{\mathrm{pro}\text{-}3}$.
At a rational coordinate cusp $P$, let $t_P$ be the standard tangent
section with leading coefficient one.  The only compact arithmetic
input is the cubic statement in Lemma~\ref{lem:cubiclocal}:
$E=J_3$ has $E(\Q)=\Z/3\Z$, $T_3\Sha(E/\Q)=0$, and
$E(\Q_3)$ is pro-$3$.  The higher Jacobian $J_n$ requires no rank
or Tate--Shafarevich hypothesis.  Instead, arithmetic localization
will kill its compact difference from a cusp.

\begin{lemma}[The inherited arithmetic action and its localization]
\label{lem:lm-power-action}
For a standard integral tangent at any rational cusp of $U_n$,
the actual action of $G_{\mathbb Q}$ on $N_n$ factors through
\[
 \Gamma=\operatorname{Gal}(K_{\{3\}}(3)/\mathbb Q),
 \qquad K=\mathbb Q(\zeta_3).
\]
Its restriction to $G_K$ has pro-$3$ image, and the local map
$G_{\mathbb Q_3}\to\Gamma$ is surjective.
If $M$ is a finite free $\mathbb Z_3$-module with an action
obtained from this action by subquotients and tensor operations,
then an unramified-outside-$3$ class
$a\in H^1(\mathbb Q,M)$ with $a_3=0$ is zero.
\end{lemma}
\begin{proof}
Write $e=3^{r-1}$.  The coordinate-power cover
$U_n\to U_3$ is geometrically Galois with group $\mu_e^2$,
of $3$-power order.  Its geometric pro-$3$ group is therefore
an open subgroup $H=N_n$ of $N_3$, and its arithmetic
pro-$3$ group is the corresponding open arithmetic subgroup.
At the cusp $(0:-1:1)$ choose the local parameters
$u_n=X/Z$ and $u_3=X/Z$ on the respective curves; the cover
satisfies $u_3=u_n^e$.  Unit tangential basepoints are thus
compatible.  The same construction applies at the other
rational cusps by coordinate permutations.

The actual upstairs action is the restriction to $H$ of the
standard downstairs action, which has already been proved to
factor through $\Gamma$ and to have pro-$3$ image over $G_K$.
The same statements follow upstairs.  The surjection
$G_{\mathbb Q_3}\to\Gamma$ is the established arithmetic
source calculation for $K$ and the single prime over $3$.
Neither claim requires taking a derived quotient of the
open embedding.

For the cohomological assertion, choose a continuous cocycle
$z$ representing $a$.  Since the action on $M$ is unramified
outside $3$, unramifiedness of the class means that this
cocycle itself is zero on all arithmetic inertia groups
outside $3$.  Its graph over $G_K$ takes values in
$M\rtimes\rho(G_K)$, a pro-$3$ group.  Consequently both
the action and the cocycle factor through $\Gamma$.
The equality $a_3=0$ permits subtracting a single coboundary
with an element of $M$ so that the new cocycle is zero on
$G_{\mathbb Q_3}$.  It still factors through $\Gamma$.
Surjectivity of the local map makes it zero globally.
\end{proof}

This also proves the compact localization assertion of
Theorem~\ref{thm:unrestrictedthreefamiliesintro} for $n=3^r$.
Indeed, a class in $\Sel(T_3J_n)$ is unramified outside three:
at every other prime the Jacobian has good reduction and its local
completed Kummer classes are unramified. Apply
Lemma~\ref{lem:lm-power-action} to the quotient $T_3J_n$ of
$N_n^{\mathrm{ab}}$. A class whose localization at three is zero
is therefore zero globally.

\begin{lemma}[Inertia of a local geometric section]
\label{lem:lm-local-inertia}
Let $p\ne3$, and let $s_p$ be the image of a full local point or cusp
section in the geometrically pro-$3$ fundamental group.  Let $J_p$ be
arithmetic inertia and let $\tau_p:J_p\to\Z_3(1)$ be the tame
$3$-adic character.  Relative to an extending tangential origin, its
geometric cocycle has the following restriction, up to geometric conjugation:
\begin{enumerate}
\item If an ordinary local point reduces into the affine special fiber,
its restriction to $J_p$ is trivial.
\item If an ordinary local point reduces to a boundary point $d$, its
restriction has image in a conjugate of cusp inertia $I_d$ and is
$\tau_p\mapsto m\tau_p$ there, where
$m\in\Z_{>0}$ is its intersection multiplicity with the boundary.
\item At a local rational cusp $d$, with parameter
$\gamma\in H^1(\Q_p,\Z_3(1))$, its restriction has image in a conjugate
of $I_d$ and is $\tau_p\mapsto v_p(\gamma)\tau_p$ there, where
$v_p(\gamma)\in\Z_3$ is the completed Kummer valuation.
\end{enumerate}
\end{lemma}
\begin{proof}
The pair $(X_n,D_n)$ has good reduction at $p$, with finite \'etale disjoint
boundary.  Prime-to-$p$ specialization shows the first assertion
\citep[Expos\'e XIII, \S2]{SGA1}; the actual pointed action is
unramified by Lemma~\ref{lem:action}.
For the second, use a boundary parameter $t$ in the strict henselian
neighborhood of the reduction of $d$.  The prime-to-$p$ geometric
fundamental group of the punctured normal disc is tame cyclic; after
passing to its pro-$3$ quotient, its covers are generated by the
Kummer covers $t^{1/3^r}$.  Pulling such a cover back along the point
gives the Kummer class of $t(x)$, whose restriction to arithmetic inertia
is $v_p(t(x))\tau_p=m\tau_p$.  This computation on all finite $3$-group
covers gives the asserted map in the full pro-$3$ group.  Changing paths
conjugates the cusp inertia subgroup.  At a cusp, the same computation
uses its Kummer parameter $\gamma$, giving the third assertion.
\end{proof}

\begin{lemma}[Compact class and boundary center]
\label{lem:lm-power-compact}
Let $s$ be a Selmer section of an arithmetic quotient of
$\pi_1(U_n)$ which retains $N_n^{\mathrm{ab}}$.
Here every local section must be the image of a full point
section or a full rational-cusp section.
There is a unique rational coordinate cusp $P\in D_n(\mathbb Q)$
with the following properties:
\begin{enumerate}
\item Every full local representative at $3$ is cuspidal at $P$.
\item The global compact class of $s$ in
$H^1(\mathbb Q,T_3J_n)$ is the Kummer class of $P$.
\item At every $p\ne3$, any local cusp representative is at $P$,
and any ordinary local representative reducing to the boundary
reduces to $P$.
\end{enumerate}
\end{lemma}
\begin{proof}
Compose the abelian section with the compact map
$X_n\to X_3=E$.  Its global class is a compact $3$-primary
Selmer class of $E$, hence is the class of one of the three
rational cubic cusps.  Injectivity of the local Kummer map
$E(\mathbb Q_3)\to H^1(\mathbb Q_3,T_3E)$ identifies the
image of the actual local center with that cusp.
Its fiber on $X_n$ consists entirely of coordinate cusps.
Because $\mathbb Q_3$ has no nontrivial roots of unity of
odd order, exactly one point of this fiber is rational over
$\mathbb Q_3$, and it is the corresponding rational cusp
$P$ of $X_n$.  Thus the local center is $P$; in particular,
the full local representative is cuspidal.

Use the standard tangent at $P$ as global origin, and let
$a\in H^1(\mathbb Q,T_3J_n)$ be the compact class of $s$
relative to that origin.  The first assertion gives $a_3=0$.
For $p\ne3$, $X_n$ and $J_n$ have good reduction.  The
compact Kummer class of every local point is unramified:
the reduction kernel in $J_n(\mathbb Q_p)$ is pro-$p$ and
uniquely $3$-divisible, and the $3$-primary completion is
the $3$-primary part of the finite reduction group.
A cusp representative has the compact Kummer class of its
center, so the same assertion holds for it.  Thus $a$ is
unramified outside $3$.  Lemma~\ref{lem:lm-power-action}
applies to $T_3J_n$, a quotient of $N_n^{\mathrm{ab}}$,
and gives $a=0$.  This proves the second assertion without
any Selmer computation for $J_n$.

Every difference of coordinate cusps on $X_n$ is killed
by $n$.  Indeed, the tangent line at any coordinate cusp
$Q$ meets $X_n$ only at $Q$, with intersection multiplicity
$n$.  For two cusps $Q,P$, the ratio of their tangent-line
equations has divisor $n(Q-P)$.  Consequently $[Q-P]$
is $3$-primary torsion in $J_n$.

At $p\ne3$, a local cusp representative centered at
$Q\in D_n(\mathbb Q_p)$ has zero compact difference from
$P$.  The $3$-primary torsion of $J_n(\mathbb Q_p)$
injects into its $3$-primary completion, so $[Q-P]=0$.
The Abel--Jacobi embedding of the positive-genus curve
$X_n$ is injective, and therefore $Q=P$.

If an ordinary local point $x$ reduces to the boundary,
its reduction lifts uniquely to
$Q\in D_n(\mathbb Q_p)$ because $D_n$ is finite \'etale
over $\mathbb Z_p$.  The difference $[x-Q]$ belongs to
the formal reduction kernel of $J_n$, so its $3$-primary
Kummer class is zero.  Since the compact class of $x$
is that of $P$, the preceding argument again gives
$[Q-P]=0$ and $Q=P$.
The initial cubic image also proves uniqueness of $P$.
\end{proof}

\begin{theorem}[Abelian Selmer cuspidality for every power of three]
\label{thm:lm-power-abelian}
For every $r\ge1$, the Selmer sections with geometric group
$N_{3^r}^{\mathrm{ab}}$ are exactly the three disjoint
rational cusp packets, each parametrized by
$H^1(\mathbb Q,\mathbb Z_3(1))$.
The statement does not assume a global full pro-$3$ lift.
Moreover, for any quotient retaining this abelianization,
every full local representative at $3$ has the same cusp
parameter as this global cusp section.  At $p\ne3$ its
whole geometric inertia is, up to geometric conjugation,
\[
 \tau\longmapsto
 i_P\bigl(v_p(\beta)\tau_p(\tau)\bigr),
 \qquad \tau\in I_p,
\]
where $\beta$ is the global parameter and $\tau_p$ is
the tame $3$-adic character.
\end{theorem}
\begin{proof}
Take the cusp $P$ of Lemma~\ref{lem:lm-power-compact}
and use its standard tangent as origin.  There is an
exact sequence of $G_{\mathbb Q}$-modules
\[
 0\longrightarrow B_n\longrightarrow N_n^{\mathrm{ab}}
 \longrightarrow T_3J_n\longrightarrow0,
 \qquad
 B_n=\mathbb Z_3(1)[D_n(\overline{\mathbb Q})]
       /\mathbb Z_3(1)\sum_{Q\in D_n}e_Q.
\]
The compact class is zero.  The equality
$(T_3J_n)^{G_{\mathbb Q}}=0$ therefore gives a unique
preimage $b\in H^1(\mathbb Q,B_n)$ of the abelian class.
These fixed points vanish because the rational torsion
of an abelian variety over a number field is finite;
the identical fixed-point vanishing holds over each
$\mathbb Q_p$ because its torsion subgroup is finite.

Choose a second rational cusp $R\ne P$, and eliminate
$e_R$ from the boundary relation.  This splits
$B_n=I_P\oplus M_n$ equivariantly over $\mathbb Z_3$.
Write $b=(\beta,b_M)$.  At $3$, local cuspidality at $P$
and the injection
$H^1(\mathbb Q_3,B_n)\hookrightarrow
 H^1(\mathbb Q_3,N_n^{\mathrm{ab}})$ give $(b_M)_3=0$.
At $p\ne3$, the geometric local inertia calculation for
the smooth pair $(X_n,D_n)$ applies.  It is zero for
an ordinary point reducing into $U_n$; at a boundary
reduction or at a cusp it lies in the inertia of its
boundary center.  Lemma~\ref{lem:lm-power-compact}
identifies that center with $P$.  Hence $b_M$ is
unramified outside $3$.

The module $M_n$ is a quotient of a submodule of
$N_n^{\mathrm{ab}}$, so its action has the properties
of Lemma~\ref{lem:lm-power-action}.  The lemma gives
$b_M=0$.  Thus the abelian section is the cusp section
$t_{P,\beta}$.

The inclusion of $I_P$ induces injections in global
and local first cohomology: it is a direct summand of
$B_n$, and the following compact Tate module has no
fixed points.  Consequently the parameter is unique,
and the full local cusp parameter at $3$ equals
$\beta_3$.  Under the power cover its inertia parameter
is multiplied by $e=3^{r-1}$; this is consistent with
$u_3=u_n^e$.  Multiplication by $e$ is injective on
$H^1(\mathbb Q_3,\mathbb Z_3(1))
 \simeq\mathbb Z_3^2$, since $\mathbb Q_3$ has no
$3$-power roots of unity.

Finally the full local inertia formula is the tame
disc calculation, with its exponent identified by
the abelian class just obtained.  The uniqueness of
the cusp center was proved in the preceding lemma.
Every global cusp section is Selmer, giving the
asserted exact description of the Selmer set.
\end{proof}

\begin{corollary}[Filling the selected cusp]
\label{cor:lm-fill-full}
For a full geometric pro-$3$ Selmer section $s$ as in
Theorem~\ref{thm:lm-power-abelian}, let $W=U_n\cup\{P\}$.
The image of $s$ in the geometrically pro-$3$ fundamental group of
$W$ is the global point section at $P$.
\end{corollary}
\begin{proof}
Filling $P$ kills the closed normal subgroup generated by $I_P$.
Consequently the inertia formula in
Theorem~\ref{thm:lm-power-abelian} makes the entire projected
cocycle unramified outside $3$, including at primes where
$v_p(\beta)\ne0$.  The action is the quotient of the action defined
by the standard tangent $t_P$, whose image in $W$ is the point
section at $P$.  Thus both its action and its cocycle factor through
$\Gamma_{\{3\}}$, as before.  At $3$, filling the local cusp section
gives the point section at $P$, so the projected local cohomology
class is trivial.  Use its local geometric conjugator globally;
the resulting cocycle still factors through $\Gamma_{\{3\}}$ and
is zero on $G_{\Q_3}$.  Surjectivity of
$G_{\Q_3}\to\Gamma_{\{3\}}$ makes it zero globally.
\end{proof}

\begin{remark}[Infinite support and central differences]\label{rem:infinite-support}
Filling the selected cusp kills its entire ramification contribution.
Before filling, a local geometric conjugator can change the full
inertia subgroup even when abelianizations agree. The lower-central
argument below controls that conjugator at each graded level; it does
not claim that subtracting the cusp parameter makes the entire cocycle
unramified.

There is no general finite-support conclusion for $\beta$.
Indeed,
\[
 H^1(\Q,\Z_3(1))\simeq
 \varprojlim_m\Q^\times/(\Q^\times)^{3^m}
\]
identifies with exponent families $(b_p)_p\in\prod_p\Z_3$ such that,
for each $m$, only finitely many $b_p$ are nonzero modulo $3^m$.
The support itself may be infinite: if $p_1,p_2,\ldots$ enumerate the
primes other than $3$, take $b_{p_j}=3^j$ and $b_3=0$.
This is an actual global cusp parameter, whose cusp section is Selmer
and ramifies at every $p_j$.  Thus continuity and local geometry cannot
by themselves force the original section into a fixed finite
ramification set.

If the parameter's support is contained in one of the finite sets for
which the arithmetic source is dominated by $G_{\Q_3}$, the inertia
formula does place the whole section in that restricted theorem and
proves cuspidality.  The unrestricted conclusion is proved below by
applying the arithmetic-source argument to successive central differences.
\end{remark}
\subsection{Free groups and free solvable central factors}
\label{subsec:lm-central-structure}

Write $N=N_n$ and
\[
 \gamma_1N=N,\qquad \gamma_{m+1}N=\overline{[N,\gamma_mN]},
 \qquad L_m=\gamma_mN/\gamma_{m+1}N.
\]
These are the closed ordinary lower-central subgroups.  The following
facts provide both the local cusp normalizer and the torsion-freeness
needed in the derived quotients.

\begin{lemma}[Free pro-$3$ facts]
\label{lem:lm-free-central}
The group $N=N_n$ is free pro-$3$ of rank $n^2+1$.  Each $L_m$ is a finite free
$\Z_3$-module, and the graded Lie algebra $\bigoplus_mL_m$ is the free
$\Z_3$-Lie algebra on $N^{\rm ab}$.  In particular there is a natural
$G_\Q$-equivariant surjection
\[
 (N^{\rm ab})^{\otimes m}\longrightarrow L_m.
\]
Moreover $\bigcap_m\gamma_mN=1$.
For any cusp $P$, its inertia $I_P$ is generated by a member of a free
pro-$3$ basis of $N$.  If $0\ne b\in\Z_3$, then
\[
 N_N(I_P^{\,b})=I_P.
\]
Here $I_P^{\,b}$ denotes the image of multiplication by $b$ in $I_P$.
\end{lemma}
\begin{proof}
The geometric presentation has $2g(X_n)$ handle generators and $3n$
inertia generators, with the single surface relation.  Eliminating any
inertia generator other than the one at $P$ leaves a free basis of size
$2g(X_n)+3n-1=n^2+1$ containing a generator $a$ of $I_P$.

For a finitely generated free pro-$3$ group, the lower central factors
are finite free $\Z_3$-modules with the Witt ranks.  This is recorded by
\citet[\S3.5, p.~71, immediately after Theorem 3.5.1]{LubotzkySegal2002}.  The commutator maps give the
surjection from the free Lie algebra; equality of the Witt ranks makes
it an isomorphism in each degree.  Iterated brackets give the displayed
tensor surjection.  Every finite $3$-group is nilpotent, so any nonidentity
element of $N$ survives in a finite quotient in which some $\gamma_mN$
vanishes.  This proves the asserted intersection.

For completeness, the normalizer assertion follows from these same Lie
facts.  An element normalizing $\langle a^b\rangle$ acts on this
procyclic group by an element of $\Z_3^\times$.  Abelianization in $N$
shows that this element is $1$, since the image of $a$ is a basis vector
and $b\ne0$.  Thus the element centralizes $a^b$.
The centralizer of the basis vector $A$ in the free Lie algebra over
$\Q_3$ is $\Q_3A$: embed the free Lie algebra in the free associative
algebra, where a homogeneous polynomial commuting with $A$ is a scalar
multiple of a power of $A$; only the first power belongs to the Lie
algebra, since every Lie element is primitive for the standard
coproduct, whereas $A^m$ is not primitive for $m>1$ in characteristic zero.  If $h\in N$ centralizes $a^b$, its degree-one class must
therefore be a $\Z_3$-multiple of $A$.  Multiply $h$ by the corresponding
power of $a$ to arrange $h\in\gamma_2N$.  Its class in each successive
$L_m$, $m\ge2$, commutes with $bA$ and hence vanishes.  The intersection
just proved gives $h=1$ after this adjustment.  Thus the original
centralizer, and therefore the normalizer, is $I_P$.
\end{proof}

\begin{lemma}[Completion of nilpotent groups with free central factors]
\label{lem:lm-nilpotent-completion}
Let $G$ be a finitely generated nilpotent group whose ordinary
lower-central factors are free abelian.  Its pro-$p$ completion
$\widehat G_p$ satisfies
\[
 \gamma_i(\widehat G_p)/\gamma_{i+1}(\widehat G_p)
 \simeq\bigl(\gamma_iG/\gamma_{i+1}G\bigr)\otimes\mathbb Z_p.
\]
\end{lemma}
\begin{proof}
Choose bases of the successive free abelian central factors,
lift them to $G$, and order the resulting Malcev basis
$a_1,\ldots,a_h$ by central weight.  Every element has a
unique collected expression $a_1^{z_1}\cdots a_h^{z_h}$
with $z_i\in\mathbb Z$.  The finite nilpotent collection
process expresses multiplication and inversion by rational
polynomials in these coordinates (equivalently by finitely
many integer-valued binomial polynomials).

The same formulas extend continuously to $\mathbb Z_p^h$.
They take values there, because they do so on the dense set
of integer tuples, and the group identities extend by
continuity.  The resulting compact group $P$ has a central
series whose factors are the corresponding powers of
$\mathbb Z_p$, so $P$ is pro-$p$.  The closure of
$\gamma_iG$ consists of the tuples whose coordinates of
weight less than $i$ vanish.  Continuity of commutators and
density give $\gamma_iP=\overline{\gamma_iG}$.

Finally, every homomorphism $G\to B$ to a finite $p$-group
is continuous for this topology.  If $p^a$ is an exponent
of $B$, its kernel contains all collected expressions
whose coordinates are divisible by $p^a$.  It therefore
contains a neighborhood of the identity and extends to $P$.
Conversely every finite continuous quotient of $P$ is a
$p$-group.  Thus $P=\widehat G_p$, and the coordinate
description of the central factors proves the formula.
\end{proof}

\begin{lemma}[Free solvable pro-$p$ central factors]
\label{lem:lm-free-solvable-pro-p-grades}
If $F$ is a free pro-$p$ group of finite rank and $F^{[d]}$
is its $d$th closed derived subgroup, then every ordinary
lower-central factor of $F/F^{[d]}$ is a finite free
$\mathbb Z_p$-module.
\end{lemma}
\begin{proof}
Let $F_0$ be the discrete free group on the same generators
and put $S_d=F_0/F_0^{[d]}$.  \citet[Theorem~3.1]{Shmelkin1964} proves that
the lower-central factors of every free polynilpotent group
are free abelian; free solvable groups are the special case
in which every entry of the polycentral class row is $2$.
Thus each nilpotent truncation
$S_d/\gamma_{c+1}S_d$ has free abelian lower-central factors.

The universal properties of free groups and pro-$p$
completion identify
\[
 \bigl(S_d/\gamma_{c+1}S_d\bigr)^{\wedge_p}
 \simeq (F/F^{[d]})/\gamma_{c+1}(F/F^{[d]}).
\]
Indeed both sides have exactly the same homomorphisms to
finite $p$-groups of derived length at most $d$ and
nilpotency class at most $c$.  Apply
Lemma~\ref{lem:lm-nilpotent-completion} to the left side.
Taking $c\ge i$ gives the assertion for the $i$th factor.
\end{proof}
The original discrete group input is
\citet[\S3, Theorem~3.1, p.~101]{Shmelkin1964}; the free-solvable
specialization is specified on p.~91.  The preceding argument proves
the passage to pro-$p$ completion.

\subsection{Central induction in the retained quotient}
\label{subsec:lm-unrestricted-full}

The induction uses full local geometric representatives supplied by
the Selmer condition.  It never requires a global lift to the full group.
The local conjugator is controlled at each central level before the
next additive difference is passed through the arithmetic source.

\begin{lemma}[Fixed points in all central quotients]
\label{lem:lm-fixed-central}
Let $R$ be a $G_\Q$-stable quotient of $N_n$ retaining its
abelianization, and suppose $L_m(R)=\gamma_mR/\gamma_{m+1}R$ is
torsion-free for every $m$. Let $t=t_{P,\beta}$ be the image in this
quotient of any global cusp section of $U_n$, with arbitrary
$\beta\in H^1(\Q,\Z_3(1))$, possibly with infinite ramification support.
For every finite prime $p$ and every $m\ge1$,
\[
 L_m(R)^{G_{\Q_p}}=0,
 \qquad
 (R/\gamma_mR)^{t(G_{\Q_p})}=1.
\]
The action on $L_m(R)$ is independent of $\beta$: it is the action defined
by the standard extending tangent $t_P$.
\end{lemma}
\begin{proof}
Inner automorphisms of $N$ act trivially on every $L_m(R)$, so changing the
section from $t_P$ to $t$ leaves this action unchanged.
For $p\ne3$, the action on $N_n^{\rm ab}$ is unramified and fits into
\[
 0\longrightarrow B_n\longrightarrow N_n^{\rm ab}
 \longrightarrow T_3J_n\longrightarrow0.
\]
Arithmetic Frobenius eigenvalues on $T_3J_n\otimes\Q_3$ have complex
absolute value $p^{1/2}$, while those on $B_n\otimes\Q_3$ have absolute
value $p$: the latter are $p$ times roots of unity from the boundary
permutation.  The tensor surjection in
Lemma~\ref{lem:lm-free-central}, followed by the surjection onto
$L_m(R)$, implies that every Frobenius eigenvalue
on $L_m(R)\otimes\Q_3$ has absolute value $p^{w/2}$ for an integer
$m\le w\le2m$.  None is $1$.  Torsion-freeness of $L_m(R)$ therefore gives
$L_m(R)^{G_{\Q_p}}=0$.

At $3$ no new $3$-adic Hodge-theoretic input is required.  The actual
action defined by $t_P$ factors through
\[
 \Gamma_{\{3\}}=\operatorname{Gal}(K_{\{3\}}(3)/\Q),
 \qquad K=\Q(\zeta_3),
\]
because it is unramified outside $3$ and its restriction to $G_K$ has
pro-$3$ image.  The same is therefore true of its action on $L_m(R)$.
The surjection in Lemma~\ref{lem:source},
$G_{\Q_3}\twoheadrightarrow\Gamma_{\{3\}}$,
gives
\[
 L_m(R)^{G_{\Q_3}}=L_m(R)^{\Gamma_{\{3\}}}
 =L_m(R)^{G_\Q}\subseteq L_m(R)^{G_{\Q_2}}=0.
\]
Finally, filter $R/\gamma_mR$ by the images of its lower central
subgroups.  A fixed element has zero image in $L_1(R)$, then in $L_2(R)$,
and successively in all $L_j(R)$ with $j<m$.  It is therefore the identity.
This applies to the possibly ramified action defined by $t$, because
its actions on the graded factors are the actions just treated.
\end{proof}

\begin{lemma}[Nonzero local inertia forces a cusp]
\label{lem:lm-nonzero-inertia-cusp}
Let $p\ne3$, and put $N=N_n$. Suppose a local section $u$ has inertia
\[
 z_u(J_p)=g I_P^{\,b}g^{-1},\qquad 0\ne b\in\Z_3,
\]
relative to $t_P$.  Then $u$ is a cusp section at $P$ in the full
geometric pro-$3$ group.  If its abelian class equals that of
$t_{P,\beta}$, it is conjugate to $(t_{P,\beta})_p$.
\end{lemma}
\begin{proof}
Arithmetic inertia is normal in $G_{\Q_p}$, and its original tangent
action on $N$ is trivial.  The cocycle identity therefore gives
\[
 z_u(\sigma\tau\sigma^{-1})
 =z_u(\sigma)\,\sigma\bigl(z_u(\tau)\bigr)\,z_u(\sigma)^{-1}
 \quad(\sigma\in G_{\Q_p},\ \tau\in J_p).
\]
Thus the image of the local section normalizes the geometric subgroup
$z_u(J_p)=g I_P^{\,b}g^{-1}$.  Conjugate by $g^{-1}$.
The original tangent action preserves $I_P$ and all its nonzero closed
subgroups.  By Lemma~\ref{lem:lm-free-central}, the normalizer of
$I_P^{\,b}$ in $N\rtimes G_{\Q_p}$ is consequently
$I_P\rtimes G_{\Q_p}$.  Thus the conjugated section takes values in
the cusp decomposition group.  Its parameter belongs to
$H^1(\Q_p,I_P)$.  The injection
$H^1(\Q_p,I_P)\hookrightarrow H^1(\Q_p,N^{\rm ab})$
proved above identifies this parameter with $\beta_p$ whenever the
abelian classes agree.
\end{proof}

\begin{theorem}[Derived and central quotients for all powers of three]
\label{thm:lm-power-all-derived}
Let $n=3^r$ with $r\ge1$, and let $R$ be a
$G_{\mathbb Q}$-stable quotient of $N_n$ retaining
$N_n^{\mathrm{ab}}$.  Suppose every ordinary lower-central
factor $\gamma_mR/\gamma_{m+1}R$ is torsion-free over
$\mathbb Z_3$.  Every Selmer section of this arithmetic
quotient is globally cuspidal.

In particular, this holds for the full group $N_n$,
for every finite lower-central quotient
$N_n/\gamma_cN_n$ with $c\ge2$, and for every finite-derived
quotient $N_n/N_n^{[d]}$ with $d\ge1$, including the maximal
geometrically metabelian pro-$3$ quotient of $U_{3^r}$.
Each Selmer set consists of exactly three disjoint
packets parametrized by $H^1(\mathbb Q,\mathbb Z_3(1))$.
\end{theorem}
\begin{proof}
Theorem~\ref{thm:lm-power-abelian} supplies the unique
abelian cusp section $t=t_{P,\beta}$.  We spell out why
the central induction now takes place directly in $R$.

Lemma~\ref{lem:lm-fixed-central} supplies
$(R/\gamma_mR)^{t(G_{\mathbb Q_p})}=1$ at every finite prime,
and its proof identifies the action on each graded factor with
the original standard-tangent action.  The finite generation and
freeness of these factors follow from
Lemma~\ref{lem:lm-free-central} and the hypothesis on $R$.

At $3$ the full local representative is conjugate
to $t$.  At $p\ne3$ with $v_p(\beta)\ne0$, the
whole inertia of a full local representative is
a nonzero subgroup of a conjugate of $I_P$.
Lemma~\ref{lem:lm-nonzero-inertia-cusp} therefore makes the
full local representative cuspidal at $P$, with parameter
$\beta_p$ by the abelian parameter injection.  At $v_p(\beta)=0$, its
whole geometric inertia is zero.  These are
statements about the full \emph{local} representatives
supplied by the Selmer condition, and may then
be projected to $R$.

Assume that $s$ and $t$ agree modulo $\gamma_mR$.
Their next difference is an additive cocycle in
$L_m(R)=\gamma_mR/\gamma_{m+1}R$.
At a place where they are locally conjugate,
the image of a local conjugator in $R/\gamma_mR$
is fixed by $t$, hence is trivial.  The conjugator
therefore lies in $\gamma_mR$.
At $3$ this proves that the graded difference
has zero local cohomology class.  At $p\ne3$ with
nonzero $v_p(\beta)$, conjugation by the cusp inertia
is inner and hence trivial on $L_m(R)$, so the
graded difference is zero on arithmetic inertia.
At the other primes both whole inertias were
already zero.  The graded difference is thus
unramified outside $3$ and zero locally at $3$.
Lemma~\ref{lem:lm-power-action} makes it a global
coboundary.  Conjugating advances the agreement
to $\gamma_{m+1}R$.

Nonempty closed sets of conjugators at the
successive levels are nested in the compact
group $R$, so they have a common element.
The closed ordinary lower-central series of a
pro-$3$ group has trivial intersection, since
every finite $3$-group is nilpotent.  The common
conjugator consequently identifies the sections
in $R$ itself.

For the listed special cases, free pro-$3$
lower-central factors are finite free.  The
same is true of free solvable pro-$3$ groups
by the discrete free-polynilpotent theorem and
the nilpotent completion argument proved in
Lemma~\ref{lem:lm-free-solvable-pro-p-grades}.
Thus all of these quotients satisfy the
hypothesis. This proves the pure-power case of
Theorem~\ref{thm:unrestrictedthreefamiliesintro}.  The proof never
asserts injectivity of a derived quotient under the open embedding
$N_n\subseteq N_3$.
\end{proof}

\begin{remark}[Why the induction is stronger than cusp subtraction]\label{rem:central-induction}
The proof never asserts that subtracting the cusp parameter makes the
whole cocycle unramified.  Instead, at each central level the local
conjugator is forced into that level by the vanishing of fixed points.
Its effect on inertia then vanishes in the next central graded module.
Only this additive difference is passed through the arithmetic quotient
$\Gamma_{\{3\}}$.  This is the step that resolves the conjugator issue
without imposing finite support on the cusp parameter.
\end{remark}

\begin{remark}[Why the proof works directly upstairs]
\label{rem:lm-derived-cover-obstruction}
An open embedding of full pro-$3$ groups need not remain
injective after taking derived quotients.  Moreover, if
$H$ is an open subgroup of a free pro-$3$ group $F$,
the quotient $F/H'$ need not have torsion-free ordinary
central factors, even though it retains $F^{\mathrm{ab}}$.
For example, let $F$ be free on $x,y$ and let $H$ be the
kernel of $F\to\mathbb Z/3\mathbb Z$ taking $x$ to $1$
and $y$ to $0$.  In $F/H'$, the relation $[x^3,y]=1$
makes $3[x,y]=0$ in the second lower-central factor.
The image of $[x,y]$ there is nonzero: map $F$ onto the
Heisenberg group over $\mathbb F_3$, where the image of
$H$ is abelian and $[x,y]$ is a nontrivial central element.
The second lower-central factor is therefore
$\mathbb Z/3\mathbb Z$.

Consequently a full-group cover argument for $U_{3^r}\to U_3$
does not by itself prove the finite-derived assertion upstairs.
Theorem~\ref{thm:lm-power-all-derived} avoids this shortcut: it first
kills the compact class by arithmetic localization, then performs
the boundary and central arguments directly in the upstairs quotient.
\end{remark}
